%%%%%%%%%%%%%%%%%%%%%%%%
%
% $Autor: Akshay Joseph $
% $Datum: 2025-04-02 $
% $Pfad: GitHub/ML24-EX01-KeywordSpotting/Nano33BLESense/Contents/General/TensorFlowLiteMirco.tex $
% $Version: 1 $
%
%%%%%%%%%%%%%%%%%%%%%%%%



%

% !TeX spellcheck = GB

% !TeX program = pdflatex

% !TeX encoding = utf8

%

%%%





\chapter{Library  \FILE{LiteRT Mirco}}

\index{LiteRTMicro}

\section{Description}



LiteRT Micro executes TensorFlow Lite models that have been converted into FlatBuffer format (\texttt{.tflite}). These models are compiled into the firmware image and executed entirely within a statically allocated memory buffer known as the tensor arena. The design avoids all dynamic memory allocation, making it deterministic and safe for use in systems with limited runtime support.



\section{Key Design Characteristics}



\begin{itemize}

	\item \textbf{No Dynamic Memory Allocation:} All required memory is statically allocated during initialization. This improves predictability and avoids heap fragmentation.

	

	\item \textbf{Header-Based Model Storage:} Models are embedded in firmware as byte arrays. This eliminates dependency on file I/O or external memory.

	

	\item \textbf{Operator Subset:} Only the operators required by the model are linked into the binary, reducing code size. Additional custom or fused operators can be added if needed.

\end{itemize}



\subsection*{Interpreter Flow}



\begin{itemize}

	\item Model is loaded into read-only memory.

	\item Tensor arena is initialized with a pre-defined buffer.

	\item Interpreter allocates tensors and invokes inference.

	\item Results are made available via output tensor access.

\end{itemize}



\section{Supported Use Cases}



LiteRT Micro is typically deployed in applications with strict size, latency, and power constraints. Examples include:



\begin{itemize}

	\item Wake-word detection (e.g., ``Hey Google'' on ARM Cortex-M4)

	\item Activity recognition using IMU data in wearable devices

	\item Environment monitoring on battery-powered sensors

	\item Predictive maintenance via vibration or current signals

	\item Simple computer vision on grayscale camera inputs (low resolution)

\end{itemize}



\section{Platform Support}



LiteRT Micro can be compiled for:



\begin{itemize}

	\item ARM Cortex-M series (e.g., M0+, M3, M4, M7, M33)

	\item RISC-V microcontrollers

	\item ESP32 and Xtensa-based MCUs

	\item Arduino-compatible platforms (e.g., Nano 33 BLE Sense)

	\item RTOS and bare-metal environments (e.g., FreeRTOS, Zephyr)

\end{itemize}



\section{Model Constraints and Deployment Workflow}



To deploy a model:



\begin{itemize}

	\item Train and export a model using TensorFlow or Keras.

	\item Convert to \texttt{.tflite} using TensorFlow Lite Converter.

	\item Apply quantization (typically to \texttt{INT8}).

	\item Convert \texttt{.tflite} to a C array (e.g., using \texttt{xxd -i}).

	\item Compile it with the embedded project, linking with LiteRT Micro.

\end{itemize}



Models suitable for LiteRT Micro are typically $\leq$250 KB in size and use fixed-point arithmetic to reduce computation and memory requirements \cite{Google:2025a}.



\section{Installation}



LiteRT Micro is not installed via a package manager. Instead, it is integrated directly into the build system of the embedded application. The runtime is part of the TensorFlow repository under \texttt{tensorflow/lite/micro}, and it is typically used in one of the following environments:



\subsection{Manual Integration (Generic Microcontrollers)}



Clone the TensorFlow repository:



\begin{lstlisting}[language=bash, caption={Cloning TensorFlow repository}]

	git clone https://github.com/tensorflow/tensorflow.git

\end{lstlisting}



Navigate to the microcontroller runtime:



\begin{lstlisting}[language=bash, caption={Navigating to LiteRT Micro directory}]

	cd tensorflow/tensorflow/lite/micro

\end{lstlisting}



Copy the relevant source files into your embedded firmware project. Alternatively, use a build system like CMake or makefiles to link the necessary components.



\subsection{Arduino Platforms (e.g., Nano 33 BLE Sense)}



\begin{itemize}

	\item Open the Arduino IDE.

	\item Install the \texttt{Arduino\_TensorFlowLite} library:

	\begin{itemize}

		\item Go to \texttt{Sketch} $\rightarrow$ \texttt{Include Library} $\rightarrow$ \texttt{Manage Libraries...}

		\item Search for \texttt{TensorFlow Lite} and install \texttt{Arduino\_TensorFlowLite}.

	\end{itemize}

	\item This provides prebuilt examples and links the LiteRT Micro runtime automatically.

\end{itemize}



\subsection{PlatformIO / Other Embedded SDKs}



\begin{itemize}

	\item Manually add the \texttt{tensorflow/lite/micro} directory into your \texttt{lib/} or \texttt{components/} folder.

	\item Ensure the include paths reference all necessary subdirectories (e.g., \texttt{micro/kernels}, \texttt{micro/memory\_planner}, etc.).

	\item A working C++ compiler (e.g., \texttt{arm-none-eabi-g++}) and toolchain for the target MCU are required.

\end{itemize}



\section{Example with Code}



Below is a minimal example showing how to use LiteRT Micro in C++ on a microcontroller. This includes model loading, interpreter setup, and inference.



\subsection{Include Necessary Headers}



\begin{lstlisting}[language=C++, caption={Including LiteRT Micro headers}]

	#include "tensorflow/lite/micro/all_ops_resolver.h"

	#include "tensorflow/lite/micro/micro_interpreter.h"

	#include "tensorflow/lite/schema/schema_generated.h"

	#include "tensorflow/lite/version.h"

\end{lstlisting}



\subsection{Load Model and Allocate Memory}



Assume the model is stored as a byte array \texttt{g\_model\_data[]}, typically generated using \texttt{xxd -i} from a \texttt{.tflite} file.



\begin{lstlisting}[language=C++, caption={Loading the model and setting up interpreter}]

	constexpr int kTensorArenaSize = 10 * 1024;

	uint8_t tensor_arena[kTensorArenaSize];

	

	const tflite::Model* model = tflite::GetModel(g_model_data);

	tflite::AllOpsResolver resolver;

	

	tflite::MicroInterpreter interpreter(model, resolver, tensor_arena, kTensorArenaSize);

	interpreter.AllocateTensors();

\end{lstlisting}



\subsection{Set Input and Run Inference}



\begin{lstlisting}[language=C++, caption={Performing inference and accessing results}]

	TfLiteTensor* input = interpreter.input(0);

	// Fill input->data.f or input->data.int8 depending on model type

	input->data.f[0] = 0.25f;  // example value

	

	if (interpreter.Invoke() != kTfLiteOk) {

		// Handle inference failure

	}

	

	TfLiteTensor* output = interpreter.output(0);

	float result = output->data.f[0];  // example for float model

\end{lstlisting}



\subsection{Notes}



\begin{itemize}

	\item For quantized models (INT8), use \texttt{input->data.int8} and cast appropriately.

	\item Ensure that the tensor arena size is tuned for your model’s memory needs.

	\item The call to \texttt{AllocateTensors()} initializes memory layout and must be called before inference.

\end{itemize}



\cite{Google:2025}







\section{Further Reading}



\nocite{Google:2025a}

\nocite{Warden:2020}



	% Überschrift ein Level unter `refsection=chapter`, also \section*:

   \printbibliography[heading=subbibliography, segment=\therefsegment]





















